% Two-frame precleaning: a talk for the NIRPS and SPIRou teams.
% Compile: pdflatex talk.tex (twice, for the frame numbers)
\documentclass[aspectratio=169,11pt]{beamer}

\usetheme{default}
\usecolortheme{seahorse}
\setbeamertemplate{navigation symbols}{}
\setbeamertemplate{footline}[frame number]
\setbeamercolor{frametitle}{fg=structure}
\setbeamerfont{frametitle}{size=\large,series=\bfseries}
\setbeamerfont{framesubtitle}{size=\small}
\setbeamertemplate{itemize item}{\color{structure}$\bullet$}
\setbeamertemplate{itemize subitem}{\color{structure}--}

\usepackage[T1]{fontenc}
\usepackage{lmodern}
\usepackage{amsmath}
\usepackage{graphicx}
\usepackage{booktabs}
\usepackage{xcolor}

\definecolor{good}{RGB}{20,120,70}
\definecolor{bad}{RGB}{176,58,46}
\definecolor{note}{RGB}{90,90,100}

\newcommand{\vtot}{V_{\mathrm{tot}}}
\newcommand{\berv}{\mathrm{BERV}}
\newcommand{\figpath}[1]{figs/#1}
\newcommand{\source}[1]{{\footnotesize\color{note}#1}}

\title{\bfseries Two-frame precleaning}
\subtitle{Separating what moves with the star from what stays with the sky,\\
before the velocities are measured}
\author{Étienne Artigau}
\institute{\texttt{github.com/eartigau/pca2d\_preclean}}
\date{\today}

\begin{document}

\begin{frame}[plain]
  \titlepage
  \begin{center}
    \small\color{note}
    NIRPS \& SPIRou teams
  \end{center}
\end{frame}

% ------------------------------------------------------------------ 2
\begin{frame}{The problem, in one sentence}
  \framesubtitle{Two things in every spectrum do not move together}

  \begin{columns}[T]
    \begin{column}{0.56\textwidth}
      \begin{itemize}
        \item The \textbf{star} is Doppler-shifted by the barycentric motion:
          tens of km/s over a year.
        \item The \textbf{sky} and the \textbf{instrument} are not. Telluric
          residuals, fringing, persistence: all fixed at the telescope.
        \item Over a campaign the star therefore \emph{slides past} the
          contamination.
      \end{itemize}

      \vspace{0.5em}
      \begin{block}{The lever}
        If a feature keeps its place on the detector while the star slides
        past it, \textbf{the feature is not the star}. One block of
        components anchored in each frame, and the two separate themselves.
      \end{block}
    \end{column}

    \begin{column}{0.42\textwidth}
      \centering
      \includegraphics[width=\textwidth]{\figpath{highpass_template}}\\
      \source{A deep line and a mask that walks across it as the BERV changes.}
    \end{column}
  \end{columns}
\end{frame}

% ------------------------------------------------------------------ 3
\begin{frame}{What it looks like in the data}
  \framesubtitle{TOI-2120, SPIRou, 316 exposures around 1267\,nm, star rest frame}
  \centering
  \vspace{-0.5em}
  \includegraphics[height=0.84\textheight]{\figpath{sequence_1267nm}}

  \vspace{-0.2em}
  \source{Each row is one exposure, ordered by BERV. Vertical structure
  belongs to the star; anything slanted does not.}
\end{frame}

% ------------------------------------------------------------------ 4
\begin{frame}{Reading that figure}
  \framesubtitle{The same seven panels, in words}

  \begin{columns}[T]
    \begin{column}{0.52\textwidth}
      \begin{enumerate}\small
        \item \textbf{The high pass.} Continuum and blaze removed. The
          \emph{slanted ridges} are the atmosphere walking across the stellar
          lines as the BERV changes through the year.
        \item \textbf{The reconstruction}: 1 star + 3 observer components,
          plus the star's mean spectrum.
        \item \textbf{Minus the observer block} -- this is what a corrected
          file holds. Scatter $0.0740 \to 0.0490$, a \textcolor{good}{34\,\%
          drop}.
        \item[5.] \textbf{Minus everything}: $0.0419$, close to photon noise.
      \end{enumerate}
    \end{column}
    \begin{column}{0.46\textwidth}
      \begin{block}{\small The key visual}
        \small
        In panel 1 the diagonals are obvious. In panel 3 they are gone, and
        \textbf{the vertical lines are untouched}. That is the whole claim of
        the method, and it is visible without any statistics.
      \end{block}

      \vspace{0.5em}
      \begin{alertblock}{\small Not a telluric model}
        \small
        $P$ and $Q$ are found by a weighted PCA \emph{on the cube itself}.
        No line list, no model atmosphere. If the fit says something sits
        still at the telescope, the campaign showed it sitting still.
      \end{alertblock}
    \end{column}
  \end{columns}
\end{frame}

% ------------------------------------------------------------------ 5
\begin{frame}{The model}
  \framesubtitle{Two blocks, one fitted amplitude each, solved per exposure}

  Let $y(x,t)$ be the high-passed log flux on a grid uniform in $\ln\lambda$:

  \begin{equation*}
    y(x,t) =
    \underbrace{\bar{S}(x-\delta x_t) + \sum_{i=1}^{n_\star} a_i(t)\,P_i(x-\delta x_t)}_{\textstyle\text{star frame}}
    \;+\;
    \underbrace{\sum_{j=1}^{n_\oplus} b_j(t)\,Q_j(x)}_{\textstyle\text{observer frame}}
    \;+\;\varepsilon
  \end{equation*}

  \vspace{0.3em}
  \begin{columns}[T]
    \begin{column}{0.49\textwidth}\small
      \begin{itemize}
        \item $\bar{S}$: the star's mean spectrum. \textbf{Coefficient fixed
          at exactly one} -- the fit can never scale the star away to explain
          something else.
        \item $P_i$: stellar variability. $Q_j$: the observer block.
        \item $a_i, b_j$: the only free parameters.
      \end{itemize}
    \end{column}
    \begin{column}{0.49\textwidth}\small
      On a log grid a Doppler shift is a pure translation:
      \begin{equation*}
        \delta x_t = \frac{\operatorname{artanh}(v_t/c)}{\Delta},
        \quad v_t = v_{\mathrm{sys}} - v_{\berv}(t)
      \end{equation*}
      \textbf{Always relativistic.} The first-order form is wrong by
      $\tfrac12 (v/c)^2 = 1.5$\,m/s at 30\,km/s, which is the whole scale of
      the measurement.
    \end{column}
  \end{columns}
\end{frame}

% ------------------------------------------------------------------ 6
\begin{frame}{How it is solved}
  \framesubtitle{Alternating, each half-step linear and closed-form}

  \begin{columns}[T]
    \begin{column}{0.54\textwidth}
      \begin{enumerate}\small
        \item \textbf{Build the cube.} Every exposure resampled onto one
          common grid with its variance and mask. Nothing co-added.
        \item \textbf{Take the star's mean} in its rest frame (weighted
          median down the time axis).
        \item \textbf{Solve the amplitudes}: with the bases fixed, each
          exposure is an independent weighted least-squares problem.
        \item \textbf{Update the bases}: weighted PCA of the residual, each
          block in its own frame.
        \item \textbf{Repeat} until $\chi^2$ stops moving.
      \end{enumerate}
    \end{column}
    \begin{column}{0.44\textwidth}\small
      \begin{equation*}
        \begin{pmatrix}\mathbf{a}\\\mathbf{b}\end{pmatrix}
        = \left(\mathbf{M}_t^{\!\top}\mathbf{W}_t\mathbf{M}_t\right)^{-1}
          \mathbf{M}_t^{\!\top}\mathbf{W}_t\left(\mathbf{y}_t - \bar{S}_t\right)
      \end{equation*}

      \vspace{0.4em}
      $\mathbf{M}_t$ stacks the shifted $P_i$ and the unshifted $Q_j$;
      $\mathbf{W}_t$ is the inverse variance.

      \vspace{0.5em}
      Writing $\bar{S}_t$ on the \emph{right-hand side} rather than as a
      column of $\mathbf{M}_t$ is exactly what pins its coefficient at one.
    \end{column}
  \end{columns}
\end{frame}

% ------------------------------------------------------------------ 7
\begin{frame}{What the components turn out to be}
  \framesubtitle{TOI-2120: rank correlations against everything in the headers}
  \centering
  \includegraphics[width=0.84\textwidth]{\figpath{correlations}}

  \vspace{0.2em}
  \begin{columns}[T]
    \begin{column}{0.32\textwidth}\footnotesize
      \textbf{Nothing follows airmass.}\\
      Largest airmass correlation anywhere: 0.18, and it is on the
      \emph{stellar} component.
    \end{column}
    \begin{column}{0.32\textwidth}\footnotesize
      \textbf{$b_3$ follows the clock} at $+0.805$ and nothing else: a
      monotonic drift over 452 days. Instrument, not weather.
    \end{column}
    \begin{column}{0.32\textwidth}\footnotesize
      \textbf{$b_2$ follows the BERV} at $-0.737$, flat in airmass. Not a
      water column: the residual of an imperfect frame separation.
    \end{column}
  \end{columns}
\end{frame}

% ------------------------------------------------------------------ 8
\begin{frame}{The high pass: a subtlety that costs m/s}
  \framesubtitle{Why the continuum filter is taken against a template}

  \begin{columns}[T]
    \begin{column}{0.46\textwidth}\small
      A Savitzky-Golay filter is a local polynomial. It does not know what a
      line is, so a deep line \textbf{drags it down}.

      \vspace{0.4em}
      Harmless if it were the same every night. It is not: LBL blanks pixels
      below $e^{-1}$ transmission, and as the BERV sweeps,
      \textbf{the hole moves across the stellar line}.

      \vspace{0.4em}
      The filter is then dragged by a different amount each exposure, leaving
      a residual that \textbf{tracks the BERV without being stellar} -- the
      single most dangerous systematic in this business.

      \vspace{0.5em}
      \textbf{Fix:} filter the residual to a median template (which has no
      holes), then add the template's own filtered version back.
    \end{column}
    \begin{column}{0.52\textwidth}
      \centering
      \includegraphics[width=\textwidth]{\figpath{highpass_template}}\\
      \source{The spectrum is held \emph{exactly} fixed; only the ten-pixel
      mask moves. Ordinary filter: 0.00230 rms. Against a template: 0.00030,
      \textcolor{good}{7.6$\times$ less}.}
    \end{column}
  \end{columns}
\end{frame}

% ------------------------------------------------------------------ 9
\begin{frame}{A warning about linearity}
  \framesubtitle{The identity that looks like it makes this pointless}

  Savitzky-Golay is linear, so one is tempted to write
  \begin{equation*}
    (s-\bar{S}) - \mathrm{SG}[s-\bar{S}] + \bar{S} - \mathrm{SG}[\bar{S}]
    \;=\; s - \mathrm{SG}[s]
  \end{equation*}

  \begin{alertblock}{That identity holds only on identical support}
    The residual $s-\bar{S}$ inherits the exposure's mask, \emph{with holes
    where the tellurics saturate}. The median template has \emph{no holes}:
    a pixel blanked tonight was measured on fifty other nights.
    \textbf{Linearity does not survive two different supports}, and the
    difference between the two sides is exactly the BERV-dependent term.
  \end{alertblock}

  \vspace{0.3em}
  \small
  Measured with a \emph{fixed} mask, the two agree to machine precision --
  which is how the identity misleads. It took holding the spectrum still and
  moving \emph{only} the mask to see the effect at all.
\end{frame}

% ------------------------------------------------------------------ 10
\begin{frame}{The headline result: a bias that follows the BERV}
  \framesubtitle{TOI-2120, velocity residual against $\vtot = v_{\mathrm{rad}} - \berv$}
  \centering
  \includegraphics[width=\textwidth]{\figpath{berv_talk}}

  \vspace{0.3em}
  \source{Near $\vtot=0$ every stellar line sits on the telluric that
  contaminates it, so anything peaking there is not stellar. Grey band:
  $|\vtot| < 4$ km/s, where the two frames nearly coincide.}
\end{frame}

% ------------------------------------------------------------------ 11
\begin{frame}{The same result, in numbers}
  \framesubtitle{A 15.8\,$\sigma$ detection becomes a non-detection}

  \centering
  \begin{tabular}{lcc}
    \toprule
     & \textbf{As delivered} & \textbf{After precleaning} \\
    \midrule
    Bias peak at $\vtot=0$ & $144.8^{+9.2}_{-8.7}$ m/s & none detected \\
    Width (FWHM)           & 4.3 km/s                  & --- \\
    Significance           & \textcolor{bad}{15.8\,$\sigma$} & \textcolor{good}{1.4\,$\sigma$} \\
    $\Delta\mathrm{BIC}$   & $+416.5$                  & $+63.5$ \\
    95\,\% upper limit     & ---                       & $|{\rm peak}| < 216.6$ m/s \\
    \midrule
    Jitter                 & 23.5 m/s                  & \textcolor{good}{11.7 m/s} \\
    \bottomrule
  \end{tabular}

  \vspace{1em}
  \begin{columns}
    \begin{column}{0.52\textwidth}\footnotesize
      \textbf{Why 4.3\,km/s is the right width:} that is about a telluric
      line convolved with the instrumental profile. If the atmosphere makes
      it, that is the width it should have.
    \end{column}
    \begin{column}{0.44\textwidth}\footnotesize
      \textbf{Caveat, stated:} amplitude and width are strongly
      anticorrelated in the posterior ($r=-0.93$), so they are not
      independently constrained.
    \end{column}
  \end{columns}
\end{frame}

% ------------------------------------------------------------------ 12
\begin{frame}{Velocities, before and after}
  \framesubtitle{TOI-2120: the same exposures, measured twice by the same LBL stage}
  \centering
  \includegraphics[height=0.70\textheight]{\figpath{lbl_velocities}}

  \vspace{0.2em}
  \small
  Nightly means, 80 shared epochs: $47.98 \to 14.20$ m/s,
  \textcolor{good}{a factor $3.38$} (paired bootstrap $3.14$--$3.65$).
  Reduced $\chi^2$: $158.7 \to 9.5$.
\end{frame}

% ------------------------------------------------------------------ 13
\begin{frame}{Where it pays, and where it does not}
  \framesubtitle{Five campaigns: the gain is ordered by SNR, not by spectral type}
  \centering
  \includegraphics[height=0.70\textheight]{\figpath{gain_vs_snr}}

  \vspace{0.1em}
  \small
  Error bars are paired bootstraps over epochs (2000 draws). Perfect rank
  ordering, but \textbf{five points and two code versions} (slide 14): a
  clear trend, not an established law.
\end{frame}

% ------------------------------------------------------------------ 14
\begin{frame}{The same five campaigns, in numbers}

  \centering
  \begin{tabular}{llrrrrr}
    \toprule
    \textbf{Star} & \textbf{Inst.} & \textbf{SNR} & \textbf{Epochs}
      & \textbf{rms in} & \textbf{rms out} & \textbf{Gain} \\
    \midrule
    TOI-2120 & SPIRou & 33  & 80  & 47.98  & 14.20 & \textcolor{good}{$3.38\times$} \\
    TOI-1452 & NIRPS  & 38  & 51  & 101.43 & 69.37 & \textcolor{good}{$1.46\times$} \\
    GJ 328   & NIRPS  & 181 & 46  & 9.34   & 9.03  & $1.03\times$ \\
    GJ 1     & NIRPS  & 233 & 140 & 2.54   & 2.56  & $0.99\times$ \\
    Proxima  & NIRPS  & 287 & 250 & 2.96   & 3.29  & \textcolor{bad}{$0.90\times$} \\
    \bottomrule
  \end{tabular}

  \vspace{1em}
  \begin{columns}[T]
    \begin{column}{0.48\textwidth}
      \begin{block}{\small The honest summary}
        \small
        \textbf{Faint targets gain; bright ones do not.} On a photon-starved
        spectrum the contamination is comparable to the noise and worth
        removing. On a bright one it is already below the noise, and fitting
        it costs more than it returns.
      \end{block}
    \end{column}
    \begin{column}{0.48\textwidth}
      \begin{alertblock}{\small Caveats}
        \small
        The four NIRPS rows \textbf{predate the fixes of 15--16 October}
        (deep line cores no longer clipped, template high pass, LBL's HE
        profile); TOI-2120 is on the current code. The table mixes two code
        versions and will be rerun on one, Proxima first.
      \end{alertblock}
    \end{column}
  \end{columns}
\end{frame}

% ------------------------------------------------------------------ 15
\begin{frame}{When not to use it}
  \framesubtitle{Failure modes found by measurement, not by argument}

  \begin{columns}[T]
    \begin{column}{0.48\textwidth}
      \begin{alertblock}{Narrow BERV span}
        The whole separation rests on the star moving relative to the sky.
        \begin{itemize}\small
          \item Wide span: \textcolor{good}{2$\times$ better}
          \item 0.7 km/s span: \textcolor{bad}{2.4$\times$ worse}
        \end{itemize}
        \small Two slices of the \emph{same} TOI-4552 campaign, identical
        otherwise. A fit with no information still returns an answer.
      \end{alertblock}

      \begin{alertblock}{Two stellar components}
        \small
        On TOI-2120 with \texttt{--n-star 2} the second amplitude starts
        following the sky, absorbs a 232 m/s shift, and the rms climbs to
        67 m/s.
      \end{alertblock}
    \end{column}
    \begin{column}{0.48\textwidth}
      \begin{alertblock}{Shared bases need shared \emph{nights}}
        \small
        Not merely similar SNR. Fitting GJ 1, Proxima and a third target on
        one basis with only two nights in common made GJ 1 \emph{worse}
        (3.04 vs 2.61 solo).
      \end{alertblock}

      \begin{block}{What it costs: samples}
        \small
        Every sample the fit could not weight is written as NaN rather than
        guessed at, and \texttt{mask: common} fixes one line set for the whole
        campaign. On TOI-2120 the files go from \textbf{27.2\,\% NaN as
        delivered to 51.3\,\%}. Worth checking on an already photon-limited
        star.
      \end{block}
    \end{column}
  \end{columns}
\end{frame}

% ------------------------------------------------------------------ 16
\begin{frame}{Practical notes for both teams}

  \begin{columns}[T]
    \begin{column}{0.49\textwidth}
      \textbf{Running it}
      \vspace{0.3em}
      \footnotesize
      \texttt{pca2d --object TOI2120 \textbackslash}\\
      \texttt{\ \ --n-star 1 --n-earth 3 \textbackslash}\\
      \texttt{\ \ --stages cube,fit,correct,figures,lbl}

      \vspace{0.6em}
      \normalsize
      \begin{itemize}\small
        \item Reader picked from \texttt{INSTRUME}, never by hand.
        \item SPIRou, NIRPS, HARPS, HARPS-N, ESPRESSO.
        \item Corrected spectra on their \emph{original} grids: any velocity
          code can read them.
        \item Every run writes its resolved config, the components, a
          coefficient CSV and a $\sim$30-page report.
      \end{itemize}
    \end{column}
    \begin{column}{0.49\textwidth}
      \textbf{Two traps worth knowing}
      \vspace{0.4em}
      \begin{alertblock}{\small LBL NIRPS\_HE + CADC}
        \small
        A one-line table bug upstream loads the \textbf{HA} instrument
        profile for HE data. Every NIRPS HE run before this was found used
        the wrong profile.
      \end{alertblock}
      \begin{alertblock}{\small Telluric fits need a template}
        \small
        Absorption must be fitted on the spectrum \emph{divided by the
        stellar template}, never raw. The tell-tale: exponents coming out
        identical with and without a template.
      \end{alertblock}
    \end{column}
  \end{columns}
\end{frame}

% ------------------------------------------------------------------ 17
\begin{frame}{Summary}

  \begin{columns}[T]
    \begin{column}{0.52\textwidth}
      \textbf{What it does}
      \begin{itemize}\small
        \item Fits a star-frame block and an observer-frame block
          \emph{simultaneously}, from the data alone.
        \item Writes back corrected spectra on their original grids.
        \item Removes a 144.8 m/s, 15.8\,$\sigma$ BERV-locked bias on
          TOI-2120, leaving a 1.4\,$\sigma$ non-detection.
      \end{itemize}

      \vspace{0.5em}
      \textbf{Where it pays}
      \begin{itemize}\small
        \item \textbf{Faint targets}: $3.4\times$ on TOI-2120,
          $1.5\times$ on TOI-1452.
        \item Bright targets: nothing, or slightly worse.
      \end{itemize}
    \end{column}
    \begin{column}{0.46\textwidth}
      \textbf{What it needs}
      \begin{itemize}\small
        \item A wide BERV span (tens of km/s).
        \item Enough epochs to build a clean median template.
        \item One stellar component, three observer ones, as the nominal
          starting point.
      \end{itemize}

      \vspace{0.8em}
      \begin{block}{\small Open questions}
        \small
        Proxima to be redone with the HARPS updates; the SNR trend rests on
        five campaigns and wants more; the NaN cost of \texttt{mask: common}
        deserves a closer look on photon-limited stars.
      \end{block}
    \end{column}
  \end{columns}

  \vspace{0.6em}
  \centering
  \small\color{note}
  \texttt{github.com/eartigau/pca2d\_preclean} \quad$\cdot$\quad
  \texttt{eartigau.github.io/pca2d\_preclean}
\end{frame}

\end{document}
